\documentclass[10pt]{amsart}
\usepackage[a4paper,margin=23mm]{geometry}
\usepackage{amsmath,amssymb,amsthm,xfrac,hyperref,xspace,graphicx}
\usepackage[english]{babel}
\usepackage[scr=rsfs,cal=euler]{mathalfa}
\newcommand{\resp}{resp.}
\newcommand{\cf}{cf.}
\newcommand{\ie}{i.e.}
\newcommand{\eg}{e.g.}
\newcommand{\loccit}{loc.\ cit.}
\newcommand{\skpt}{\par\smallskip\noindent}
\let\Subsection\subsection
\let\Subsubsection\subsubsection
\emergencystretch=3em
\pagestyle{plain}
\multlinegap0pt

\hyphenation{inclu-de orthog-o-nal}

\renewcommand{\theenumi}{\roman{enumi}}

\newenvironment{smallpmatrix}{\bigl(\begin{smallmatrix}}{\end{smallmatrix}\bigr)}

\newcommand\mto{\mathchoice{\longmapsto}{\mapsto}{\mapsto}{\mapsto}}

\newcommand{\Psfrac}[2]{(\sfrac{#1}{#2})}
\newcommand{\psfrac}[2]{\sfrac{(#1)}{#2}}
\newcommand{\spfrac}[2]{\sfrac{#1}{(#2)}}
\newcommand{\pspfrac}[2]{\sfrac{(#1)}{(#2)}}

\let\sep:

\newcommand{\ext}{\mathrm{ext}}
\newcommand{\res}{\mathrm{res}}

\DeclareMathOperator{\dL}{\mathit{dL}}
\DeclareMathOperator{\Id}{Id}

\newtheorem{theorem}{Theorem}[section]
\newtheorem{cor}[theorem]{Corollary}
\newtheorem{lemme}[theorem]{Lemma}
\newtheorem{prop}[theorem]{Proposition}

\theoremstyle{definition}
\newtheorem{df}[theorem]{Definition}
\newtheorem{remark}[theorem]{Remark}

\let\eps\varepsilon
\let\om\omega
\let\phi\varphi
\let\e e
\let\EE\exp
\let\dis\relax
\let\ov\overline
\let\wt\widetilde
\let\<\langle
\let\>\rangle

\newcommand\C{\mathbb C}
\newcommand\R{\mathbb R}
\newcommand\T{\mathbb T}
\newcommand\Z{\mathbb Z}
\newcommand\HH{\mathbb H}
\newcommand\N{\mathbb N}
\newcommand\E{\mathcal E}
\newcommand\ET{\widetilde{\mathcal{E}}}
\newcommand\W{\mathcal W}
\newcommand\M{\mathcal M}
\newcommand\TT{\mathcal T}
\renewcommand{\L}{\mathscr{L}}
\newcommand{\dd}{\mathrm{d}}

\newcommand{\pa}[1]{\left(#1\right)}
\newcommand{\bpa}[1]{\bigl(#1\bigr)}
\newcommand{\Bpa}[1]{\Bigl(#1\Bigr)}
\newcommand{\bbpa}[1]{\biggl(#1\biggr)}
\newcommand{\pac}[1]{\left[#1\right]}
\newcommand{\bpac}[1]{\bigl[#1\bigr]}
\newcommand{\Bpac}[1]{\Bigl[#1\Bigr]}
\newcommand{\bbpac}[1]{\biggl[#1\biggr]}
\newcommand{\paa}[1]{\left\{#1\right\}}
\newcommand{\bra}[1]{\left\langle #1 \right\rangle}
\newcommand{\abs}[1]{\left|#1\right|}
\newcommand{\Babs}[1]{\Bigl|#1\Bigr|}
\newcommand{\bbabs}[1]{\biggl|#1\biggr|}
\newcommand{\absabs}[1]{\left\|#1\right\|}
\newcommand{\Babsabs}[1]{\Bigl\|#1\Bigr\|}
\newcommand{\bbabsabs}[1]{\biggl\|#1\biggr\|}

\DeclareMathOperator{\Ree}{{\mathfrak{Re}}}\let\Re\Ree
\DeclareMathOperator{\Imm}{{\mathfrak{Im}}}\let\Im\Imm
\DeclareMathOperator{\Sp}{Sp}
\newcommand{\simx}[2]{\underset{#1\to#2}{\thicksim}}

\renewcommand{\enlargethispage}[1]{}
\begin{document}
\title{Exponential linear instability of every rectangular one-flux lattice for the lowest Landau level equation under strip-periodic perturbations}
\author{Pierre Germain \and Valentin Schwinte \and Laurent Thomann}
\maketitle
\section{Lowest Landau equation and lattice definitions}
\subsection{The equation}

Consider the lowest Landau level equation\vspace*{-3pt}
\begin{equation}\label{LLL}\tag{LLL}
\begin{cases}
i \partial_t u = \Pi (|u|^2 u),\\
u(0,\cdot)= u_0,
\end{cases}
\end{equation}
where $\Pi$ is the projector in $L^2(\C)$ on the Fock-Bargmann space\vspace*{-3pt}
\[
\mathcal{E} =\bigl\{ u(z) = \EE\bigl(-{|z|^2}/2\bigr) f(z),\;f \; \text{entire\ holomorphic}\bigr\}\cap L^2(\C),
\]
which is given by the kernel\vspace*{-3pt}
\begin{equation}\label{defPi}
[\Pi u](z) = \frac{1}{\pi} \EE\bigl(-{|z|^2}/2\bigr) \int_\mathbb{C} \EE\bigl(\ov w z - |w|^2/2\bigr) u(w) \dL(w),
\end{equation}
($L$ being simply the Lebesgue measure on $\C$). Notice that $\Pi $ extends as an operator from $\mathscr{S}'(\C)$ onto the space\vspace*{-3pt}
\[
\widetilde {\mathcal E}=\bigl\{ u(z) =\EE\bigl(-{|z|^2}/2\bigr) f(z),\;f \; \text{entire\ holomorphic}\sep\exists M,\, \vert u(z)\vert \lesssim \< z\>^M \, \bigr\},
\]
on which $\Pi $ is the identity operator. For this reason, we~shall extend equation (\ref{LLL}) to $\widetilde {\mathcal E}$.

\subsection{The theta function and the Abrikosov lattice} We now turn to the mathematical description of the Abrikosov lattice, which relies on the Jacobi theta functions.

Let $\tau=\tau_1+i \tau_2 \in \C$ with $\tau_2>0$. The Jacobi theta function on the lattice $ \mathcal{L}_\tau = \mathbb{Z}\oplus \tau \mathbb{Z}$ is defined (see \cite[Chap.\,V]{Cha}) by
\begin{equation*}
\Theta_{\tau}(z)=-i \sum_{n=-\infty}^{+\infty} (-1)^n \EE\Bigl({i\pi \tau (n+1/2)^2}+{i(2n+1)\pi z}\Bigr), \qquad z\in \C.
\end{equation*}
The $\Theta_{\tau}$ function vanishes exactly on $\mathcal{L}_\tau$ and is such that
\begin{equation}\label{star}
\Theta_{\tau}(-z)=-\Theta_{\tau}(z),\quad \Theta_{\tau}(z+1)=-\Theta_{\tau}(z),\quad \Theta_{\tau}(z+\tau)=-e^{-i \pi \tau}e^{-i 2\pi z}\Theta_{\tau}(z).
\end{equation}
For more results on theta functions, we~refer to \cite[\S 4]{Godement} (in this book, the notation $\theta_1(z, \tau)=\Theta_{\tau}(z)$ is used) and we refer as well to the book \cite{Mumford}. We~also mention the expository paper \cite{Nier2} of F.\,Nier making connections with various point of views.

To every theta function, we~associate the function
\[
\Phi_0(z) = \EE\Bigl(\frac{z^2}2 - \frac{{i \pi}z} {\gamma} - \frac{|z|^2}2\Bigr) \,\Theta_\tau \Bigl(\frac{z}{\gamma} - \frac{\tau-1}{2} \Bigr).
\]
For the choice of $\gamma^2 \tau_2 = \gamma^2 \Im \tau= \pi$, this function is $L^\infty$ and belongs to $\widetilde {\mathcal E}$. Furthermore, it is a stationary solution of the \eqref{LLL} equation (see Remark \ref{rem14} below). Finally, its period is the lattice $\mathcal{L}_{\tau,\gamma}=\gamma (\mathbb{Z} + \tau \mathbb{Z})$;
if $\tau = j = \EE\bigl({{2i \pi}/{3}}\bigr)$ and $\gamma^2 = {2\pi}/{\sqrt{3}}$, this period is given the hexagonal array known as the Abrikosov lattice!

Introduce now the \textit{magnetic translation} in the direction $\alpha \in \mathbb{C}$:
\[
R_\alpha u (z) = \EE\Bigl(\frac{\alpha \ov{z} - \ov{\alpha} z}2\Bigr) u(z+\alpha).
\]
For any $\alpha\in \C$, $R_\alpha$ is a symmetry of \eqref{LLL} (and in particular, it leaves $\mathcal{E}$ invariant). The symmetries of the $\Theta_{\tau}$ function translate into the following identity for $\Phi_0$:
\begin{equation}
\label{symPhi}
R_{\gamma} \Phi_0 = R_{\gamma \tau} \Phi_0 = \Phi_0.
\end{equation}

\section{(LLL) on doubly periodic domains (cells)}
\label{section2}

In this section, we~construct solutions to \eqref{LLL} such that $\vert u\vert$ is doubly periodic. We~will rely on results of Aftalion-Serfaty~\cite[\S 3]{AftaSerfa}. In the sequel, we~work with the general lattice
\[
\mathcal{L}_{\tau,\gamma} = \gamma(\mathbb{Z}\oplus \tau \mathbb{Z})
\]
with $\gamma>0 $, and we consider the space
\begin{align*}
\mathcal{E}_{\tau,\gamma}&=\bigl\{u\in \widetilde{\mathcal{E}}: R_\gamma u=u,\ R_{\gamma \tau} u=u\bigr\}\\
&=\Bigl\{u\in \widetilde{\mathcal{E}}: u(z+\gamma)=\EE\Bigl({\frac{\gamma}2(z-\ov{z})}\Bigr) u(z), \\
&\qquad u(z+\gamma\tau)=\EE\Bigl(\frac{\gamma}2 \bigl(\tau_1(z-\ov{z})-i\tau_2(z+\ov{z})\bigr)\Bigr) u(z)=\EE\Bigl(\frac{\gamma}{2}(\ov{\tau}z-\tau\ov{z})\Bigr) u(z)\Bigr\},
\end{align*}
so that in particular for all $u \in \mathcal{E}_{\tau,\gamma}$ and $z\in\mathbb{C}$, $\vert u(z+\gamma)\vert= \vert u(z+\gamma\tau)\vert=\vert u(z)\vert$. 

We define the fundamental cell of $\mathcal{L}_{\tau,\gamma}$ by
\begin{equation*}
K_{\tau,\gamma}=\bigl\{z=\gamma(r_{1}+r_{2}\tau),\ r_{1},r_{2}\in[0,1]\bigr\}.
\end{equation*}

\Subsection{Multiplicative description of $\E_{\tau,\gamma}$}

\begin{prop}[\cite{AftaSerfa}]\label{prop41}
Let $\gamma>0$ and $\tau \in \C$ with $\tau_2=\Im \tau >0$.
\begin{enumerate}
\item If $\gamma^2 \tau_2 \notin\pi \mathbb{N}$, then $\mathcal{E}_{\tau,\gamma} =\{ 0 \}$.
\item If $\gamma^2 \tau_2 = \pi N$ for some $N \in \mathbb{N}$, then $\mathcal{E}_{\tau,\gamma}$ is a complex vector space of dimension~$N$. It~can be described as the set of functions vanishing $N$ times modulo~$\mathcal{L}_{\tau,\gamma}$ (counting multiplicity), which can be written under the form
\begin{equation}\label{formprop}
u(z)=\lambda \EE\bigl(z^2/2 + b z- |z|^2/2\bigr) \prod_{j=1}^N\Theta_{\tau}\bigl(({z-z_j})/{\gamma} \bigr),
\end{equation}
where $\lambda \in \C$ and $(z_j)_{1\leq k\leq N}\in \C$ are representatives of the zeroes of $u$ modulo~$\mathcal{L}_{\tau,\gamma}$ and finally $b \in \mathbb{C}$ and the $(z_j)$ satisfy the relations
\begin{align}\label{somme}
\gamma b &= i(- N + 2k) \pi, \nonumber \\
\sum_{j=1}^N z_j &= \frac{\gamma}{2} (\tau - 1) N - k \tau \gamma + \ell \gamma,
\end{align}
for some $k,\ell \in \mathbb{Z}$.
\end{enumerate}
\end{prop}
\paragraph{The lattice function.}
\begin{equation}\label{defPhi0}
\Phi_{0}(z)=\EE\Bigl({\frac{1}{2} z^2 -\frac{i\pi}{\gamma} z - \frac{1}{2} |z|^2} \Bigr) \Theta_{{\tau}/{N}}\bigl((z-z_{0})/{\gamma}\bigr), \quad z_0=\frac{\gamma}2 \Bigl(\frac{\tau}N-1 \Bigr),
\end{equation}
\subsection{The nonlinear term}

The aim of this subsection is to understand the structure of the nonlinear term $\Pi \bigl(|u|^2 u\bigr)$ in the space $\mathcal{E}_{\tau,\gamma}$, more specifically in the basis provided by the $\Phi_k$. The first step is to realize that $\Pi$ can be interpreted as an orthogonal projector, as will now be explained. Define the space\vspace*{-3pt}
\begin{multline*}
\mathcal{F}_{\tau,\gamma}=\Bigl\{u\in L^2 (K_{\tau,\gamma}), u(z+\gamma)=\EE\Bigl({\frac{\gamma}2(z-\ov{z})}\Bigr) u(z),\\
u(z+\gamma\tau) =\EE\Bigl(\frac{\gamma}2(\ov{\tau}z-\tau\ov{z})\Bigr) u(z)\Bigr\},
\end{multline*}
so that\vspace*{-3pt}
\[
\mathcal{E}_{\tau,\gamma} = \mathcal{F}_{\tau,\gamma} \cap \widetilde{\mathcal{E}}.
\]
\begin{lemme}
\label{lemmaproj}
For all $u,v\in \mathcal{F}_{\tau,\gamma}$,\vspace*{-3pt}
\begin{equation*}
\int_{K_{\tau,\gamma}} \Pi u(z) \ov{v(z)}\dL(z)= \int_{K_{\tau,\gamma}} u(z) \ov{\Pi(v(z))}\dL(z).
\end{equation*}
\end{lemme}

\begin{proof} With the change of variable $\xi=w-k\tau-\ell$ and the fact that\vspace*{-3pt}
\[
u(\xi+k\tau \gamma+\ell \gamma)=\EE\Bigl({\ell \gamma(\xi-\ov\xi)/2+k\gamma(\ov{\tau}\xi-\tau\ov\xi)/2}\Bigr)u(\xi),
\]
we get\vspace*{-3pt}
\begin{align*}
\Pi u(z)&=\frac1{\pi}\EE\bigl({-|z|^{2}/2}\bigr)\int_{\C}\EE\bigl(z\ov{\om}-|w|^{2}/2\bigr)\,u(w) \dL(w)\\
&=\frac1{\pi}\EE\bigl({-|z|^{2}/2}\bigr)\sum_{k,\ell\in \Z}\int_{w\in K_{\tau,\gamma}+k\tau \gamma+\ell \gamma}\EE\bigl(z\ov{\om}-|w|^{2}/2\bigr)\,u(w) \dL(w)\\
&=\frac1{\pi}\int_{ K_{\tau,\gamma}}\sum_{k,\ell\in \Z}\EE\bigl({-|z|^{2}/2}\bigr) \EE\Bigl({-{\gamma^2}|k\tau+\ell|^{2}/2+(k\ov\tau+\ell)\gamma z}\Bigr)\\
&\hspace{4cm}\times \EE\Bigl({(z-k\gamma \tau-\ell\gamma)\ov{\om}-|w|^{2}/2}\Bigr)\,u(w) \dL(w).
\end{align*}
Then\vspace*{-3pt}
\begin{multline*}
\int_{K_{\tau,\gamma}} \Pi u(z) \ov{v(z)}\dL(z)=\\
\begin{aligned}
&=\frac1{\pi}\int_{ K_{\tau,\gamma}\times K_{\tau,\gamma}}\sum_{k,\ell\in \Z}\EE\bigl({-|z|^{2}/2}\bigr) \EE\Bigl({-{\gamma^2}|k\tau+\ell|^{2}/2+(k\ov\tau+\ell)\gamma z}\Bigr)\\
&\hspace{3.5cm}\times \EE\Bigl({\bigl(z-k\gamma \tau-\ell\gamma\bigr)\ov{\om}-|w|^{2}/2}\Bigr)u(\om)\ov{v(z)}\dL(w) \dL(z)\\
&=\int_{K_{\tau,\gamma}} u(w) \ov{\Pi(v(w))}\dL(w)
\end{aligned}
\end{multline*}
hence the result.
\end{proof}

Given an element $u$ of $\mathcal{F}_{\tau,\gamma}$, it can be thought of as a function on $\mathbb{C}$, or as a function on $K_{\tau,\gamma}$. Therefore we introduce the following notations.

\begin{df}
Given an element $u$ of $\mathcal{F}_{\tau,\gamma}$, we~write $u^{\ext}$ (extended), if $u$ is considered as a function on $\mathbb{C}$ and $u^{\res}$ (restricted) if $u$ is considered as a function on~$K_{\tau,\gamma}$.
\end{df}

With these notations, we~have the following result:

\begin{lemme} \label{lemmaproj2}
The projector $\Pi$ on $\mathcal{F}_{\tau,\gamma}$ can be interpreted as the orthogonal projector~$\Pi'$ in $L^2(K_{\tau,\gamma})$ on $\mathcal{E}_{\tau,\gamma}$. In other words,
\[
(\Pi u^{\ext})\big|_{K_{\tau,\gamma}} = \Pi' u^{\res}.
\]
\end{lemme}

\begin{proof} On the one hand, if $u \in \mathcal{E}_{\tau,\gamma}$, then both $\Pi$ and $\Pi'$ act as the identity. It~is clear as far as $\Pi'$ is concerned. Turning to $\Pi$, we~have that $\Pi u^{\ext} \in \widetilde{\mathcal{E}}$ and $\Pi u^{\ext} \in \mathcal{F}_{\tau,\gamma}$ (since $R_\alpha$ and $\Pi$ commute), hence $\Pi u$ belongs to $\widetilde{\mathcal{E}} \cap \mathcal{F}_{\tau,\gamma} = \mathcal{E}_{\tau,\gamma}$. Furthermore, by the previous lemma, $\langle \Pi u^{\ext}, v \rangle_{L^2(K_{\tau,\gamma})} = \langle u^{\res}, v \rangle_{L^2(K_{\tau,\gamma})}$ for any $v \in \mathcal{E}_{\tau,\gamma}$, hence $u = \Pi u$.

On the other hand, choose $u$ in the orthogonal set to $\mathcal{E}_{\tau,\gamma}$ in $L^2(K_{\tau,\gamma})$. By~definition of $\Pi'$, we~get $\Pi' u^{\res} = 0$. Turning to $\Pi u^{\ext}$, it is in $\mathcal{E}_{\tau,\gamma}$ by the same argument as above, and satisfies $\langle \Pi u^{\ext}, v \rangle_{L^2(K_{\tau,\gamma})} = \langle u^{\res}, v \rangle_{L^2(K_{\tau,\gamma})}=0$ for any $v \in \mathcal{E}_{\tau,\gamma}$; hence $\Pi u^{\ext} = 0$.
\end{proof}

By the previous lemma, we~will identify henceforth $\Pi$ and $\Pi'$, and we will also ignore the distinction between $u^{\ext}$ and $u^{\res}$.
\skpt\begin{remark}\label{rem14}
\begin{itemize}
\item
When $N=1$, we~obtain the stationary solution
\begin{equation*}
u(t,z)= e^{i\mu t}\Phi_0(z) =e^{i\mu t}\EE\Bigl({ \frac{z^2}2 -\frac{i\pi}{\gamma} z - \frac{|z|^2}2}\Bigr) \Theta_{\tau}\Bigl(\frac{1}{\gamma}(z-z_0)\Bigr), \quad z_0 = \frac{\gamma}{2}(\tau-1).
\end{equation*}
This stationary solution has already been obtained in \cite{ABN}. Indeed, set $\gamma=1/{\nu}$, then the function $f_\tau$ which is defined in \cite[Th.\,1.4]{ABN} is the function
\[
v(z)=\EE\bigl(z^2/2-|z|^2/2\bigr)\Theta_\tau({z}/{\gamma}).
\]
Set $u:=R_{-z_0}v$, where $z_0 = {\gamma}(\tau-1)/2$, then $u \in \mathcal{E}_{\tau,\gamma}$ with $\tau_2=\pi/\gamma^2$. Notice however that $v \notin \mathcal{E}_{\tau,\gamma}$ because $ R_\gamma v=-v$ and $ R_{\gamma \tau} v=-v$.
\end{itemize}
\end{remark}
\section{(LLL) on simply periodic domains (strips)}
\label{section3}

\subsection{Periodic functions in the Fock-Bargmann space}
Denote by $\mathcal{F}_{\gamma}$ the space
\begin{equation*}
\mathcal{F}_{\gamma}=\bigl\{u\in \widetilde{{\mathcal{E}}}\sep R_{\gamma}u=u\bigr\}.
\end{equation*}
A subspace of $\mathcal{F}_\gamma$ is given by functions which are $L^p$-integrable on a fundamental domain associated to the translation $z \mto z+\gamma$. For convenience, we~fix the vertical strip
\[
S_\gamma = \bigl\{ z \in \mathbb{C}, \, -{\gamma}/{2} \leq \Re z \leq {\gamma}/{2} \bigr\},
\]
and we consider the natural $L^p(S_\gamma)$ norm. Similarly, $L^{p,\alpha}(S_\gamma) \cap \mathcal{F}_{\gamma} $ is the subspace of~$\mathcal{F}_\gamma$ for which the following norm is finite
\[
\| u \|_{L^{p,\alpha}(S_\gamma)} = \| \langle z \rangle^\alpha u \|_{L^p(S_\gamma)}.
\]
Denote
\[
\psi_0(z):=\left(\frac{2}{\pi \gamma^2}\right)^{1/4}\EE\Bigl({\frac{z^2}2-\frac{|z|^2}2}\Bigr),
\]
and for $n\in \Z$
\begin{equation*}
\psi_n(z) = R_{{i n \pi}/ \gamma}\psi_0(z) = \EE\Bigl({-\frac{\pi^2 n^2}{\gamma^2}}\Bigr) \EE\bigl({{2in\pi}{\gamma} z}\bigr)\psi_0(z).
\end{equation*}

\begin{remark}
Alternatively, we~could have defined
\[
\wt{\psi}_n(z) =R_{{n\tau \gamma}/ N}\psi_0(z)
\]
with the usual quantification condition on $\tau$ and $N$. This results in
\[
\wt{\psi}_n(z) =R_{{n\tau \gamma}/ N}\psi_0(z)= \EE\Bigl(\frac{i n^2 \pi \tau_1}{N} -\frac{\pi^2 n^2}{\gamma^2}+ \frac{2in\pi}{\gamma} z\Bigr) \psi_0(z).
\]
So the only dependence on $\tau$ is through the number $\EE\bigl({i {n^2 \pi \tau_1}/ {N}}\bigr)$ which has absolute value $1$ and is thus irrelevant.
\end{remark}

\Subsection{On the restriction of the projector $\Pi$ on $\mathcal{F}_\gamma$}

\begin{lemme}\label{lemme_extension} Let $\Pi$ be the orthogonal projector from the space of $R_\gamma$-invariant functions to ${L^2(S_\gamma) \cap \mathcal{F}_{\gamma}}$, and denote $K(z,w)$ its kernel, for $(z,w) \in S_\gamma$.
\begin{enumerate}
\item\label{lemme_extensioni} The kernel $K$ is given by the formula
\begin{multline}\label{perio}
K(z,w)=\sqrt{\frac{2}{\pi \gamma^2}} \EE\Bigl({-\frac{\pi^2}{2\gamma^2}} \Bigr)\EE\Bigl({\frac{1}{2} z^2 - \frac{1}{2} |z|^2+\frac{1}{2} \ov{w}^2 - \frac{1}{2} |w|^2 -\frac{i\pi}{\gamma} (z- \ov{w})}\Bigr)\\
\times \Theta_{{2i\pi}/{\gamma^2}}\Bigl(\frac{1}{\gamma}(z-\ov{w}-\frac{i\pi}{\gamma}+\frac{\gamma}{2})\Bigr).
\end{multline}
\item\label{lemme_extensionii} It enjoys the bound
\begin{equation}\label{borneK}
|K(z, w)| \lesssim \EE\Bigl(-\frac{(\Im z-\Im w)^2}2\Bigr).
\end{equation}
\item\label{lemme_extensioniii} For all $u,v\in \mathcal{F}_{\gamma} \cap L^2(S_{\gamma})$
\begin{equation}\label{aa}
\int_{S_{\gamma}} \Pi u(z) \ov{v(z)}\dL(z)= \int_{S_{\gamma}} u(z) \ov{\Pi(v(z))}\dL(z).
\end{equation}
\item\label{lemme_extensioniv} The projector $\Pi$ on $\{ u, \; R_\gamma u = u \} \cap L^2(S_\gamma)$ can be identified with the orthogonal projector from $L^2(S_\gamma)$ to $\mathcal{F}_\gamma \cap L^2(S_\gamma)$.
\item\label{lemme_extensionv} For any $p \geq 2$ and $\alpha \geq 0$, the orthogonal projector $\Pi$ has a unique extension to $L^p(S_\gamma)$ and $L^{p,\alpha}(S_\gamma)$.
\end{enumerate}
\end{lemme}

\skpt
\begin{proof}
\eqref{lemme_extensioni} We show that the periodized projection operator in the whole space agrees with the projection operator given by \eqref{perio}. To get a formula for the periodized projection operator, we~write, if $u$ is $R_\gamma$-invariant,
\[
u(z+k\gamma) = \EE\bigl({{k\gamma}(z-\ov{z})/2 }\bigr) u(z).
\]
The change of variable $w = w'+ k \gamma$ results in
\[
\EE\Bigl({\ov{w} z - \frac{1}{2}|w|^2}\Bigr) u (w) = \EE\Bigl(\overline{w'} z - \frac{1}{2} |w'|^2 + k \gamma (z-\overline{w'}) - \frac{1}{2} k^2 \gamma^2\Bigr) u(w').
\]
Making this change of variable in the integral after splitting it into vertical strips yields
\begin{align*}
\Pi u(z) & = \frac{1}{\pi} \EE\bigl(-{|z|^2}/2\bigr) \int_\mathbb{C} \EE\bigl(\ov w z - |w|^2/2\bigr) u(w) \dL(w) \\
&= \frac{1}{\pi} \EE\bigl(-{|z|^2}/2\bigr) \sum_{ k \in \Z} \int_{w \in S_{\gamma} + k\gamma} \dots \\
& = \frac{1}{\pi} \EE\bigl(-{|z|^2}/2\bigr) \int_{S_{\gamma}} \EE\bigl({\ov w z - |w|^2/2 }\bigr) \\[-5pt]
&\hspace{4cm} \times \biggl[ \sum_{ k \in \Z} \EE\bigl({k \gamma (z-\overline{w}) - k^2 \gamma^2/2}\bigr) \biggr] u(w) \dL(w).
\end{align*}
Comparing with the above formula, we~need to check that
\begin{multline*}
\sqrt{\frac{2}{\pi \gamma^2}} \EE\Bigl({z^2/2 +\ov{w}^2/2 }\Bigr)\sum_{ k \in \Z } \EE\Bigl({\frac{2ik\pi}{\gamma}(z-\ov{w}) -\frac{2\pi^2k^2}{\gamma^2}}\Bigr)\\
= \frac{1}{\pi} e^{\ov w z } \sum_{k \in \Z} \EE\Bigl({k \gamma (z-\overline{w}) - \frac{1}{2} k^2 \gamma^2}\Bigr),
\end{multline*}
but this follows from the Poisson formula \eqref{PF1} with $x={(z-\ov{w})}/ \gamma$ and $\alpha={\gamma^2}/2$. This proves \eqref{perio}.

\eqref{lemme_extensionii}
Let $N=2$, $\tau = {2i\pi}/{\gamma^2}$, and $z_1 = z_2 = {\gamma}(\tau-1)/2$. Then, the function
\[
u(z) = \EE\Bigl({z^2/2-{2i\pi}z/ {\gamma}-|z|^2/2}\Bigr) \Theta^2_{{2i\pi}/{\gamma^2}}\bigl((z-z_1)/\gamma\bigr) 
\]
is of the form \eqref{formprop}. Therefore, $u$ is periodic over the lattice $\mathcal{L}_{\tau,\gamma}$, so that it is bounded. This implies that
\[
\Babs{\Theta_{{2i\pi}/{\gamma^2}}\Bpa{\frac{1}{\gamma}(z-\ov{w})-\frac{i\pi}{\gamma}+\frac{\gamma}{2})}\EE \Bigl({-\frac{i\pi}{\gamma}(z-\ov{w})+\frac14(z-\ov{w})^2-\frac14|z-\ov{w}|^2} \Bigr)} \lesssim 1.
\]
Then, using \eqref{lemme_extensioni} and
\[
\frac14(z-\ov{w})^2-\frac14|z-\ov{w}|^2 = \frac14z^2-\frac14|z|^2+\frac14\ov{w}^2-\frac14|w|^2-\frac14(2z\ov{w}-zw-\ov{z}\ov{w}), 
\]
we get
\[
\abs{K(z,w)} \lesssim \Babs{\EE \Bigl(\frac14z^2-\frac14|z|^2+\frac14\ov{w}^2-\frac14|w|^2-\frac14(2z\ov{w}-zw-\ov{z}\ov{w}) \Bigr)}.
\]
We write, for $z=x+iy, w=a+ib \in S_{\gamma}$,
\begin{multline*}
\frac14z^2-\frac14|z|^2+\frac14\ov{w}^2-\frac14|w|^2-\frac14(2z\ov{w}-zw-\ov{z}\ov{w})\\
\begin{aligned}
& = \frac{i}{2}xy-\frac12y^2-\frac{i}{2}ab-\frac{1}{2}b^2 + yb +2i(ya-xb) \\
& = -\frac12(y-b)^2 + \frac{i}2\bigl(xy-ab+2(ya-xb)\bigr),
\end{aligned}
\end{multline*}
and obtain the bound
\[
\abs{K(z,w)} \lesssim \EE\Bigl(-\frac{(\Im z-\Im w)^2}2\Bigr).
\]

\eqref{lemme_extensioniii} Can be proved like Lemma \ref{lemmaproj}.

\eqref{lemme_extensioniv} Can be proved like Lemma \ref{lemmaproj2}.

\eqref{lemme_extensionv} By the assertion \eqref{lemme_extensionii}, we~can write for $u\in L^p_{\gamma}$
\begin{align*}
\|\Pi u\|_{L^{p,\alpha}(S_{\gamma})} &= \bbabsabs{\bra{z}^\alpha\int_{S_{\gamma}}K(z,w)u(w)\dL(w)}_{L^p(S_{\gamma})}\\
& \lesssim \bbabsabs{\int_{S_{\gamma}} \EE\Bigl(-\frac{(\Im z-\Im w)^2}2\Bigr) \bigl(\bra{w}^\alpha+\bra{z-w}^\alpha\bigr)u(w)\dL(w)}_{L^p(S_{\gamma})} \\
& \lesssim \absabs{\bra{z}^\alpha u}_{L^p(S_{\gamma})} = \|u\|_{L^{p,\alpha}(S_{\gamma})},
\end{align*}
and the result follows.
\end{proof}

We shall now prove hypercontractivity estimates for all $L^p-L^q$.
\begin{lemme}
Let $u\in \mathcal{F}_{\gamma}$. Then for any $1 \leq p \leq q \leq +\infty$,
\begin{equation}\label{Carlen_bande}
\|u\|_{L^{q}(S_{\gamma})} \lesssim \|u\|_{L^{p}(S_{\gamma})}.
\end{equation}
\end{lemme}

\begin{proof}
By interpolation, it is sufficient to prove the result for $p \geq 1$ and $q= +\infty$. If $u\in \mathcal{F}_\gamma$,
\[
u(z)= \Pi u(z)=\int_{S_\gamma} K(z,w) u(w) \dL(w).
\]
By \eqref{borneK}, the kernel $K$ and the Hölder inequality, for all $z \in S_{\gamma}$,
\[
|u(z)|\leq \| K(z, \cdot)\|_{L^{p'}(S_{\gamma})} \| u\|_{L^p(S_{\gamma})} \lesssim \|u\|_{L^p(S_{\gamma})},
\]
which is the desired result.
\end{proof}

\Subsection{The family $(\psi_n)_{n \in \Z}$ is a Hilbertian basis of $L^2(S_\gamma) \cap \mathcal{F}_\gamma$}

\begin{lemme}\label{lem98}
The family $(\psi_n)_{n \in \Z}$ is orthonormal and forms a Hilbertian basis of $L^2(S_\gamma) \cap \mathcal{F}_\gamma$.
\end{lemme}

In particular, the function $\Phi_0$ defined in \eqref{defPhi0} has a natural expansion in the family $(\psi_n)_{n \in \Z}$, namely a direct computation shows that, with $z_0={\gamma}({\tau}/N-1)/2$,
\begin{align}\label{exp1}
\Phi_{0}(z)&=\EE\Bigl(\frac{z^2}2 -\frac{i\pi z} {\gamma} - \frac{|z|^2}2\Bigr) \Theta_{{\tau}/{N}}\bigl((z-z_{0})/\gamma\bigr) \nonumber \\
&= e^{-i{\pi \tau}/{(4N)}}\pa{{\pi\gamma^2}/{2}}^{1/4}\sum_{n=-\infty}^{+\infty} e^{i{n^2\pi\tau_1}/{N}}\psi_n(z).
\end{align}

\begin{proof} It is clear that $\psi_0 \in \mathcal{F}_\gamma$, and this remains true for $\psi_n$ since $\big[R_{{i n \pi}/\gamma},R_\gamma\big] = 0$. Furthermore,
\begin{multline*}
\int_{S_{\gamma}}\hspace*{-1mm} \psi_{k}(z)\ov{\psi_{\ell}} \dL(z) =\left(\frac{2}{\pi \gamma^2}\right)^{1/2}\hspace*{-1mm}
\int_{-\infty}^\infty \int_0^\gamma \hspace*{-1mm}\EE\Bigl(-2y^2+\frac{2i\pi}{\gamma}(k-\ell)x-\frac{2\pi}{\gamma}(k+\ell)y\Bigr) dx dy \\
= \left(\frac{2}{\pi \gamma^2}\right)^{1/2} \delta_{k,\ell} \gamma \int_{-\infty}^\infty \EE\Bigl(-2y^2-\frac{4\pi}{\gamma}ky\Bigr) dy = \delta_{k,\ell},
\end{multline*}
which shows that the family is orthonormal. 

We now show that the family is complete. Let $u \in L^2(S_\gamma) \cap \mathcal{F}_\gamma$. For $z=x+iy$, we~set $v(x,y)=u(z)\EE\bigl(-ixy\bigr)$. Then for all $y \in \R$, the function $x \mto v(x,y)$ is $\gamma$-periodic. We~expand it in Fourier series $\dis v(x,y)=\sum_{k \in \Z} c_k(y)\EE\bigl({{2i \pi k x}/\gamma}\bigr)$ and we have
\[
\|u \|^2_{L^2(S_\gamma)}=\|v\|^2_{L^2(S_\gamma)}=\sum_{k \in \Z} \int_{\R}|c_k(y)|^2 dy <\infty. 
\]
Therefore, for all $k \in \Z$, $c_k \in L^2(\R)$. We~now expand $c_k$ in the Hilbertian basis of~$L^2(\R)$ given by (translated) Hermite functions. Namely, for all $k\in \Z$, there exists a family of polynomial $(P_{j,k})_{j \geq 0}$ where $P_{j,k}$ has degree $j$ and such that
\[
c_k(y)=\EE\Bigl(-\bigl(y+\frac{k \pi}\gamma\bigr)^2\Bigr)\sum_{j=0}^{+\infty}P_{j,k}(y).
\]
Coming back to $u$, it can be written as the following convergent series
\[
u(z)=\sum_{k \in \Z} \sum_{j\geq 0} \EE\Bigl({ixy+\frac{2i \pi k x}\gamma -\bigl(y+\frac{k \pi}\gamma\bigr)^2} \Bigr) P_{j,k}(y),
\]
where the summands are orthogonal. To show that the family $(\psi_k)$ is complete, using the continuity of $\Pi$ on $L^2(S_\gamma)$ by Lemma \ref{lemme_extension}\eqref{lemme_extensioniii}, it suffices to show that the orthogonal projection of each summand can be written as a linear combination of these functions.

We will do so by showing that, for all $k \in \Z$ and $j \geq 0$, there exists $C_{j,k} \in \C$ such that
\begin{equation}\label{claim1}
A_{j,k}=\Pi\Bigl(\EE\Bigl({-\bigl(y+\frac{k \pi}\gamma\bigr)^2}\Bigr) \bigl(y+\frac{k \pi}\gamma\bigr)^j \EE\Bigl({ixy+\frac{2ik\pi x}{\gamma} } \Bigr)\Bigr)= C_{j,k} \psi_k(z),
\end{equation}
By the formula for the projection
\begin{align*}
A_{j,k}
&= \frac{1}{\pi} e^{-|z|^2/2} \int_\mathbb{C} \EE\bigl({\ov w z - |w|^2/2}\bigr) \EE\Bigl({-\bigl(\frac{w -\ov w}{2i}+\frac{k \pi}\gamma\bigr)^2} \Bigr)(\frac{w -\ov w}{2i}+\frac{k \pi}\gamma)^j \\
&\hspace{5.5cm}\times \EE\Bigl({\frac14(w^2-\ov w^2) +\frac{i \pi k}{\gamma}(w+\ov w)} \Bigr) \dL(w)\\
&= \frac{1}{\pi (2i)^j} e^{-|z|^2/2} \int_\mathbb{C} \EE \Bigl({-|w|^2 + \ov w z +\frac{2i \pi k}{\gamma} w+\frac12 w^2 } \Bigr) \bigl({w -\ov w}+\frac{2i k \pi}\gamma\bigr)^j \dL(w)\\
&= \frac{2^{j/2+1}}{\pi (2i)^j} e^{-|z|^2/2} \int_\mathbb{C} \EE\Bigl({-2|w|^2 + \sqrt{2}\ov w z +\frac{2 \sqrt{2}i \pi k}{\gamma} w+ w^2 } \Bigr) \\
&\hspace{5.5cm}\times \bigl({w -\ov w}+\frac{\sqrt{2}i k \pi}\gamma\bigr)^j \dL(w).
\end{align*}
Let $t\in \R$ and let us compute
\begin{multline*}
\sum_{j=0}^{+\infty} \frac{\pi (2i)^j}{2^{j/2+1}}\frac{t^j A_{j,k}}{j!}=\EE\Bigl(-\frac{|z|^2}2 + \frac{\sqrt{2}i k \pi t } \gamma \Bigr)\\
\times \int_\mathbb{C} \EE\Bigl({-2|w|^2 +\bigl(\frac{2 \sqrt{2}i \pi k}{\gamma} +t\bigr)w + (\sqrt{2}z-t)\ov w + w^2 } \Bigr) \dL(w).
\end{multline*}
From \cite[\S 6]{GGT} we recall the formula
\begin{equation*}
\int \EE\Bigl(-2 |w|^2 + aw + b\overline{w} + cw^2 \Bigr)\dL(w)
= \frac{\pi}{2} \EE\bigl({b(cb+2a)}/{4}\bigr).
\end{equation*}
We set $a={2 \sqrt{2}i \pi k}/{\gamma} +t$, $b=\sqrt{2}z-t$ and $c=1$, thus
\begin{multline}\label{series}
\sum_{j=0}^{+\infty} \frac{\pi (2i)^j}{2^{j/2+1}}\frac{t^j A_{j,k}}{j!}=\frac{\pi}{2} \EE\Bigl(-\frac{|z|^2}2 + \frac{\sqrt{2}i k \pi t } \gamma \Bigr)\\
\times \int_\mathbb{C} \EE\Bigl({-2|w|^2 +\bigl(\frac{2 \sqrt{2}i \pi k}{\gamma} +t\bigr)w + (\sqrt{2}z-t)\ov w + w^2 } \Bigr) \dL(w) \\
= \frac{\pi}{2} e^{- t^2/4} \EE\Bigl({\frac{z^2}{2}-\frac{|z|^2}2+\frac{2i k\pi z}{\gamma}}\Bigr) = C_k e^{-t^2/4}\psi_k(z).
\end{multline}
Finally, by identifying the powers of $t$ in the expansion \eqref{series} we get \eqref{claim1}.
\end{proof}

\begin{remark}

Since the $(\psi_k)$ form a Hilbertian basis of $L^2(S_\gamma) \cap \mathcal{F}_\gamma$, we~can recover the expression of the kernel $K$ as follows:
\begin{align*}
K(z,w)&= \sum_{ k \in \Z } \psi_k(z)\ov{\psi_k(w)}\\
&= \sqrt{\frac{2}{\pi \gamma^2}} \EE\Bigl({\frac12z^2-\frac12|z|^2+\frac12\ov{w}^2-\frac12|w|^2}\Bigr) \sum_{ k \in \Z } \EE\Bigl({\frac{2ik\pi}{\gamma}(z-\ov{w}) -\frac{2\pi^2k^2}{\gamma^2}}\Bigr)\\
&=\sqrt{\frac{2}{\pi \gamma^2}} \e^{-{\pi^2}/{(2\gamma^2)}} \Theta_{{2i\pi}/{\gamma^2}}\Bigl(\frac{1}{\gamma}(z-\ov{w}-\frac{i\pi}{\gamma}+\frac{\gamma}{2})\Bigr) \\
&\hspace{3.5cm} \times \EE\Bigl({-\frac{i\pi}{\gamma}(z-\ov{w}) +\frac12z^2-\frac12|z|^2+\frac12\ov{w}^2-\frac12|w|^2}\Bigr).
\end{align*}
\end{remark}
\subsection{Cubic coefficients on the strip}
From now on, we~denote
\begin{equation}\label{def-A}
A_{k,\ell,m} = \frac{1}{\gamma\sqrt{\pi}}\EE\Bigl(-\frac{\pi^2}{\gamma^2}\bigl((\ell-k)^2+(\ell-m)^2\bigr)\Bigr).
\end{equation}
Then, the following result holds true.

\begin{lemme}\label{lem-ortho}
If $k,\ell,m,n \in \mathbb{Z}$ are such that $k - \ell + m - n \neq 0$, then
\begin{equation*}
\int_{S_{\gamma}} \psi_{k}\ov{\psi_{\ell}} \psi_{m}\ov{\psi_{n}}(z)\dL(z)
=\begin{cases}
\displaystyle A_{k,\ell,m} & \text{if $k - \ell + m - n = 0$}, \\
0 & \text{if $k - \ell + m - n \neq 0$}.
\end{cases}
\end{equation*}

In other words, when $m = k - \ell + n$, we~have
\begin{equation}\label{proj3}
\Pi\bigl(\psi_{k}\ov{\psi_{\ell}} \psi_{m}\bigr)
=A_{k,\ell,m} \psi_{n}.
\end{equation}
\end{lemme}

\begin{proof} We write\vspace*{-3pt}
\begin{multline*}
\int_{S_{\gamma}} \psi_{k}\ov{\psi_{\ell}} \psi_{m}\ov{\psi_{n}}(z)\dL(z)
= \frac{2}{\pi \gamma^2}\int_{0}^{\gamma}\int_{\R}
\EE\Bigl(\frac{2i \pi}{\gamma}(k-\ell+m-n)x -\frac{2\pi}{\gamma}(k+\ell+m+n)y\Bigr)\\
\times \EE\Bigl(-4y^2-\frac{\pi^2}{\gamma^2}(k^2+\ell^2+m^2+n^2)\Bigr) dy \, dx,
\end{multline*}
which vanishes if $k-\ell+m-n \neq 0$. Assuming that $k-\ell+m-n =0$,\vspace*{-3pt}
\begin{multline*}
\int_{S_{\gamma}} \psi_{k}\ov{\psi_{\ell}} \psi_{m}\ov{\psi_{n}}(z)\dL(z)=\\
\begin{aligned}
&= \frac{2}{\pi \gamma} \int_{\R} \EE\Bigl(-\frac{2\pi}{\gamma}(k+\ell+m+n)y-4y^2-\frac{\pi^2}{\gamma^2}(k^2+\ell^2+m^2+n^2)\Bigr) dy\\
&= \frac{2}{\pi \gamma}\EE\Bigl(\frac{\pi^2}{4 \gamma^2}(k+\ell+m+n)^2-\frac{\pi^2}{\gamma^2}(k^2+\ell^2+m^2+n^2)\Bigr) \int_{\R} \e^{ -4y^2} dy,
\end{aligned}
\end{multline*}
which gives the desired formula.
\end{proof}

Expanding a solution of~\eqref{LLL} in the Hilbertian basis
\[
u(t,z) = \sum_{k \in \mathbb{Z}} \lambda_k(t) \psi_k(z),
\]
it follows from the above proposition that the coordinates $(\lambda_n)$ satisfy the equation
\begin{equation}\label{eq_lambda_k}
i\frac{d}{dt} \lambda_n(t) = \sum_{k - \ell + m = n} A_{k,\ell,m} \lambda_{k}(t) \ov{\lambda_{\ell}(t)} \lambda_{m}(t), \quad n \in \Z, \quad t\in \R.
\end{equation}
\setcounter{section}{3}
\section{Linearized analysis around rectangular lattices}\label{rectangular_lattice}
\label{section4}

From now on and in the following sections, we~assume that $N=1$. In the present section, we~assume furthermore that the lattice is rectangular, namely that $\tau_1=0$ and $\tau_2={ \pi}/{\gamma^2}$, so that $\tau= { i\pi}/{\gamma^2}$.

Recall the definition of the strip
\[
S_{\gamma}=\bigl\{z=x+iy\sep x\in[0,\gamma],\; y\in \R\bigr\},
\]
which we regard as a fundamental domain for functions in $\mathcal{F}_\gamma$.

Also recall the Hilbert basis
\begin{align*}
\psi_0(z)&=\Bigl(\frac{2}{\pi \gamma^2}\Bigr)^{1/4}\EE\Bigl(\frac{z^2}2-\frac{|z|^2}2\Bigr),\\
\psi_k(z)&=R_{{ik\pi}/{\gamma}}\psi_0(z)=\Bigl(\frac{2}{\pi \gamma^2}\Bigr)^{1/4}\EE\Bigl(\frac{2ik\pi} {\gamma} z -\frac{\pi^2k^2}{\gamma^2}+\frac{z^2}{2}-\frac{|z|^2}2\Bigr), \quad k \in \mathbb{Z}.
\end{align*}
It is such that $R_{\alpha}\psi_k=\EE\bigl({{2ik\pi \alpha}/ {\gamma}}\bigr) \psi_k$ if $\alpha \in \mathbb{R}$ and $\Gamma_1 \psi_k=\bigl({2 \pi k}/{\gamma} \bigr)\psi_k$.

By \eqref{exp1} we have the expansion
\begin{align*}
\Phi_0(z)&= \Bigl(\frac{\pi \gamma^2}{2}\Bigr)^{1/4} \EE\Bigl(\frac {\pi^2}{4\gamma^2}\Bigr) \sum_{n=-\infty}^{+\infty} \psi_n (z)\\
&=\EE\Bigl(\frac{z^2}2 -\frac{i\pi z} {\gamma} - \frac{|z|^2}2\Bigr) \Theta_{{i\pi}/{\gamma^2}}\Bigl({ \bigl({z-\frac{{i\pi}/{\gamma}-\gamma}2}}\bigr)/\gamma\Bigr),
\end{align*}
and it was observed in Remark \ref{rem14}, that this function gives the $M$-stationary wave
\[
\Phi(t,z) = e^{-i\lambda t}\Phi_0(z).
\]
Ignoring its phase, this solution is periodic with respect to translations in $\gamma \mathbb{Z}$ and~$\gamma \tau \mathbb{Z}$, and thus corresponds physically to a rectangular lattice.
The aim of this section is to linearize \eqref{LLL} around the rectangular lattice, and describe the spectrum of this linearization.

Linearizing \eqref{LLL} equation around $\Phi_0$ yields
\begin{equation}
i\partial_t u = \Pi\big[2|\Phi|^2u+\Phi^2\ov{u}\big] = \Pi\bpac{2|\Phi_0|^2u+e^{-2i\lambda t}\Phi_0^2\ov{u}},
\end{equation}
so that the function $v = \EE\bigl({i\lambda t}\bigr)u$ satisfies the equation
\begin{equation}\label{Lin_LLL}
i\partial_t v + \lambda v = \Pi\bpac{2|\Phi_0|^2v+\Phi_0^2\ov{v}}.
\end{equation}

\begin{theorem}
\label{thmunstable}
Writing
\[
v(t) = \sum_{n \in \Z} c_n(t) \psi_n \qquad \text{and} \qquad
f(t,\xi) = \sum_{n \in \Z} c_n(t) e^{-2\pi i n\xi},
\]
Equation \eqref{Lin_LLL} becomes
\[
i \partial_t \begin{pmatrix} {f}(t,\xi) \\[2pt] \overline{{f}}(t,-\xi) \end{pmatrix} = A_{\mathrm{rect}}(\xi) \begin{pmatrix} {f}(t,\xi) \\[2pt] \overline{{f}}(t,-\xi) \end{pmatrix},
\]
for a $2\times 2$ matrix $A(\xi)$. The matrix $A_{\mathrm{rect}}(\xi) $ can be diagonalized for $\xi \neq 0$, resulting~in
\[
A_{\mathrm{rect}}(\xi) = P(\xi) \begin{pmatrix} \mu(\xi) & 0 \\[2pt] 0 & -\mu(\xi) \end{pmatrix} P^{-1}(\xi)
\]
(we refer to the proof for exact formulas for $A$, $P$ and $\mu$, which are somewhat lengthy and therefore omitted here).

Furthermore,
\begin{itemize}
\item for any $\gamma>0$, $\mu(\xi) \in i \mathbb{R} \setminus\{ 0 \}$ if $\xi$ is close to $0$;
\item if $\tau = i$ (square lattice), then $\mu(\xi) \in i \mathbb{R} \setminus\{ 0 \}$ for any $\xi \neq 0$.
\end{itemize}

In particular, all rectangular lattices are exponentially unstable in $L^2(S_\gamma) \cap \mathcal{F}_\gamma$.
\end{theorem}
\begin{proof}[Proof. The linearized operator in the Hilbert basis $(\psi_n)$]
Setting\vspace*{-5pt}\enlargethispage{\baselineskip}
\[
\dis v(t,z)=\sum_{n=-\infty}^{+\infty} c_n(t)\psi_n(z),
\]
the equation \eqref{Lin_LLL} becomes\vspace*{-5pt}
\begin{equation*}
i\partial_t c_n + \lambda c_n = \pa{\frac{\pi}{2}}^{1/2}\gamma \EE\Bigl(\frac{\pi^2}{2\gamma^2}\Bigr)\sum_{\substack {k, \ell, m \in \Z\\ k-\ell+m=n}}A_{k,\ell,m}(2c_{k} + \ov{c_{\ell}}),
\end{equation*}
thanks to the equation \eqref{proj3}, recalling the notation \eqref{def-A}\vspace*{-3pt}
\[
A_{k,\ell,m} = \frac{1}{\gamma\sqrt{\pi}}\EE\Bigl({-\frac{\pi^2}{\gamma^2}\pa{(\ell-k)^2+(\ell-m)^2}}\Bigr).
\]
For $n\in \Z$,\vspace*{-5pt}
\begin{align*}
\sum_{\substack {k, \ell, m\in \Z\\ k-\ell+m=n}}A_{k,\ell,m}c_{k} &= \frac{1}{\gamma\sqrt{\pi}}\sum_{k \in \Z} \EE\Bigl(-\frac{\pi^2}{\gamma^2}(n-k)^2\Bigr)c_{k} \sum_{\substack {\ell,m \in \Z\\ m-\ell=n-k}} \EE\Bigl(-\frac{\pi^2}{\gamma^2}(\ell-k)^2\Bigr) \\[-3pt]
&=\frac{1}{\gamma\sqrt{\pi}}\sum_{k \in \Z} \EE\Bigl(-\frac{\pi^2}{\gamma^2}(n-k)^2\Bigr)c_{k}\sum_{q\in\Z}\EE\Bigl(-\frac{\pi^2q^2}{\gamma^2}\Bigr),
\end{align*}
and\vspace*{-5pt}
\begin{align*}
\sum_{\substack {k, \ell, m\in \Z\\ k-\ell+n=m}}A_{k,\ell,m}\ov{c_{\ell}} &= \frac{1}{\gamma\sqrt{\pi}}\sum_{\ell \in \Z} \ov{c_{\ell}} \sum_{\substack {k,n \in \Z\\ k+m=n+\ell}}\EE\Bigl({-\frac{\pi^2}{\gamma^2}\bigl((n-k)^2+(\ell-k)^2\bigr)}\Bigr) \\[-3pt]
&= \frac{1}{\gamma\sqrt{\pi}}\sum_{\ell \in \Z} \ov{c_{\ell}} \sum_{k\in\Z}\EE\Bigl({-\frac{\pi^2k^2}{\gamma^2}}\Bigr)\EE\Bigl({-\frac{\pi^2}{\gamma^2}(n-\ell-k)^2}\Bigr). \\
\end{align*}
In other words,\vspace*{-5pt}
\begin{equation*}
i\partial_t \begin{pmatrix} c_n \\[2pt] d_n \end{pmatrix} = \begin{pmatrix} \mathcal{L}^{\mathrm{rect}}-\lambda \Id & \mathcal{M}^{\mathrm{rect}} \\[2pt] -\mathcal{M}^{\mathrm{rect}} & - (\mathcal{L}^{\mathrm{rect}}-\lambda \Id) \end{pmatrix} \begin{pmatrix} c_n \\[2pt] d_n \end{pmatrix},
\end{equation*}
where $(d_n):= (\ov{c_n})$ and\vspace*{-3pt}
\begin{gather}
\mathcal{L}^{\mathrm{rect}}: u \mto C_L L * u, \quad C_L =\sqrt{2}\EE\Bigl({\frac{\pi^2}{2\gamma^2}}\Bigr)\sum_{q\in\Z}\EE\Bigl(-\frac{\pi^2q^2}{\gamma^2}\Bigr),\notag
\\[-5pt]
\label{def_op_L_rect}
L_n = \EE\Bigl(-\frac{\pi^2n^2}{\gamma^2}\Bigr),
\\
\mathcal{M}^{\mathrm{rect}}: u \mto C_M M * u, \quad C_M =\frac{1}{\sqrt{2}}\EE\Bigl({\frac{\pi^2}{2\gamma^2}}\Bigr),\notag
\\[-5pt]
\label{def_op_M_rect}
M_n = \sum_{p=-\infty}^{+\infty}\EE\Bigl({-\frac{\pi^2p^2}{\gamma^2}}\Bigr) \EE\Bigl({-\frac{\pi^2}{\gamma^2}(n-p)^2}\Bigr) = [L * L](n),
\end{gather} 
where $*$ stands for the discrete convolution of sequences.

\subsubsection*{The linearized operator in the Fourier variable $\xi$}
Define the discrete Fourier transform $\dis \mathcal{F}:\ell^2(\Z) \to L^2([0,1])$(respectively the inverse Fourier transform $\mathcal{F}^{-1}$) for a sequence $\dis u = (u_n)_{n\in\Z} \in \ell^2(\Z)$(respectively a function $f \in L^2([0,1])$) by
\begin{equation}
\label{deffourier}
\mathcal{F}(u)(\xi) = \sum_{k= -\infty}^{+\infty} u_ke^{- 2i\pi k \xi}, \qquad \mathcal{F}^{-1}(f)(k) = \int_0^1 e^{2i\pi k \xi} f(\xi) d\xi.
\end{equation}
As is well-known, $\mathcal{F}(u*v) = \mathcal{F}(u) \mathcal{F}(v)$, so that
\begin{align*}
\mathcal{F}\mathcal{L}^{\mathrm{rect}}\mathcal{F}^{-1}f&=C_L\mathcal{F}(L) f, \\
\mathcal{F}\mathcal{M}^{\mathrm{rect}}\mathcal{F}^{-1}f&=C_M\mathcal{F}(M) f = C_M\mathcal{F}(L*L) f = C_M \mathcal{F}(L)^2 f.
\end{align*}
Turning to the Fourier transform of $L$, is is given by
\begin{equation}\label{def_fct_l}
\begin{aligned}
\mathcal{F}(L)(\xi) = \sum_{k= -\infty}^{+\infty} L_ke^{-{2i\pi k}\xi} &= \sum_{k = -\infty}^{+\infty} \EE\Bigl({-{2i\pi k}\xi-\frac{\pi^2 k^2}{\gamma^2}}\Bigr)\\
&=1 + 2 \sum_{k=1}^{+\infty}\EE\Bigl({-\frac{\pi^2k^2}{\gamma^2}}\Bigr)\cos\pa{2\pi k\xi}=: \ell(\xi).
\end{aligned}
\end{equation}

All functions and constants above can now be expressed through $\ell$ and $C_M$:
\[
C_L = 2 C_M \ell(0), \quad \lambda = C_M \ell^2(0), \quad \mathcal{F}(M)(\xi) = \ell^2(\xi).
\]
Overall, we~find the expression
\begin{equation*}
\pac{\mathcal{F}\begin{pmatrix} \mathcal{L}^{\mathrm{rect}}-\lambda \Id & \mathcal{M}^{\mathrm{rect}} \\[2pt] -\mathcal{M}^{\mathrm{rect}} & - (\mathcal{L}^{\mathrm{rect}}-\lambda \Id) \end{pmatrix}\mathcal{F}^{-1}(X)}\pa{\xi} =:A_{\mathrm{rect}}(\xi)X,
\end{equation*}
where
\begin{equation*}
A_{\mathrm{rect}}(\xi) = \begin{pmatrix} a(\xi) & b(\xi) \\[2pt] -b(\xi) & -a(\xi) \end{pmatrix} =C_M \begin{pmatrix} \ell(0)(2\ell(\xi)-\ell(0)) & \ell^2(\xi) \\[2pt] -\ell^2(\xi) & -\ell(0)(2\ell(\xi)-\ell(0)) \end{pmatrix}.
\end{equation*}

\subsubsection*{Diagonalizing the matrix $A_{\mathrm{rect}}$}
The characteristic polynomial of $A_{\mathrm{rect}}(\xi)$ is
\[
P_{A_{\mathrm{rect}}(\xi)}(X) = X^2 + D(\xi), \quad \text{where} \quad D(\xi) = \det(A_{\mathrm{rect}}(\xi)) = b^2(\xi) - a^2(\xi).
\]
Whether the eigenvalues at frequency $\xi$ are real or imaginary (and therefore, stable or unstable) depends on the sign of $D(\xi)$, which can be factorized as follows
\begin{align*}
D(\xi) &= C_M^2\pac{\ell^2(\xi)+\ell(0)(2\ell(\xi)-\ell(0))}\pac{\ell^2(\xi)-\ell(0)(2\ell(\xi)-\ell(0))} \\
&= \underbrace{C_M^2\bpac{\ell(\xi)-\ell(0)}^2\bpac{\ell(\xi)+(1+\sqrt{2})\ell(0)}}_{\geq 0}\bpac{\ell(\xi)-(\sqrt{2}-1)\ell(0)}.
\end{align*}
(where the above term is non-negative since $|\ell(\xi)|$ is maximum at $\xi=0$). We~will denote
\[
\mu(\xi) = \sqrt{-D(\xi)},
\]
with the convention that $\Im  \sqrt x \geq 0$ if $x \leq 0$. The eigenvalues of $A_{\mathrm{rect}}(\xi)$ are $\mu_{\pm} (\xi)= \pm \sqrt{-D(\xi)}$, and the corresponding eigenvectors $e_{\pm}(\xi) = (b(\xi), \mu_{\pm}(\xi) - a(\xi))$. Denoting~$P$ for the change of base matrix $P(\xi) = (e_+(\xi) | e_-(\xi))$, we~obtain the diagonalization formula
\[
A_{\mathrm{rect}}(\xi) = P(\xi) \begin{pmatrix} \mu(\xi) & 0 \\ 0 & -\mu(\xi) \end{pmatrix} P^{-1}(\xi).
\]
\subsubsection*{The sign of the determinant}
As demonstrated above, $D(\xi)$ and
\[
F(\xi):=\ell(\xi) -(\sqrt{2}-1)\ell(0)
\]
have the same sign. By~continuity of $F$, it appears that $F(\xi)$, and hence $D(\xi)$, is~positive in a neighborhood of $\xi=0$.

To obtain a more precise bound, we~resort to the geometric series formula to bound~$F$ from below:
\begin{align*}
F(\xi) &=2 - \sqrt{2} + 2 \sum_{k=1}^{+\infty}\EE\Bigl(-\frac{\pi^2 k^2}{\gamma^2}\Bigr) \bpa{\cos\pa{2\pi k\xi}-\sqrt{2}+1}\\
&\geq 2 - \sqrt{2} - 2 \sum_{k=1}^{+\infty}\sqrt{2}\EE\Bigl(-\frac{\pi^2 k^2}{\gamma^2}\Bigr) \geq 2 - \sqrt{2} - 2\sqrt{2}\,\frac{q}{1-q},
\end{align*}
where $q:= \EE\bigl(-{\pi^2}/{\gamma^2}\bigr)$. Then, $F(\xi) > 0$ as soon as
\[
q < \frac{2-\sqrt{2}}{2\sqrt{2}+2-\sqrt{2}} \approx 0.1716,
\]
which is the case for the square lattice, for which $\gamma^2 = \pi$, so that we have \linebreak ${q = \EE\bigl({-\pi} \bigr)\approx 0.0432\dots}$

\subsubsection*{Linear instability} We claim that the equation \eqref{Lin_LLL} is unstable in $L^2(S_\gamma) \cap \mathcal{F}_\gamma$. We~note first that
\[
\| v(t,z) \|_{L^2(S_\gamma)} = \| c_n(t) \|_{\ell^2(\mathbb{Z})} = \| f(t,\xi) \|_{L^2([0,1])}
\]
(where the first equality is a consequence of the fact that $(\psi_n)$ is a Hilbert basis, and the second equality is Parseval's theorem).

As a consequence, it suffices to establish instability in the $ {f}$ variable. Set $f_0(\xi)=f(0,\xi)$. By~the above diagonalization formula, $ {f}$ satisfies
\[
\begin{pmatrix} {f}(t,\xi) \\[2pt] \overline{{f}}(t,-\xi) \end{pmatrix} = P(\xi) \begin{pmatrix} e^{-i t \mu(\xi)} & 0 \\[2pt] 0 & e^{i t \mu(\xi)}\end{pmatrix} P^{-1}(\xi) \begin{pmatrix} {f}_0(\xi) \\[2pt] \overline{{f}_0}(-\xi) \end{pmatrix}.
\]
As we saw above, $\Re (-i \mu(\xi)) > 0$ for $\xi$ close to zero. Therefore, exponential growth will be observed as soon as the first coordinate of $P^{-1}(\xi) \begin{smallpmatrix} {f}_0(\xi) \\\overline{{f}_0}(-\xi) \end{smallpmatrix}$ is non-zero. This is easy to arrange, and it is even the generic situation! As an example, we~give $v(0) = \psi_0$, leading to $ {f}_0(\xi) = 1$, and
\[
P^{-1}(\xi) \begin{pmatrix} {f}_0(\xi) \\[2pt] \overline{{f}_0}(-\xi) \end{pmatrix} = \begin{pmatrix}\frac{b(\xi)+a(\xi)}{\mu(\xi)(b(\xi)+a(\xi)-\mu(\xi))} \\[2pt] \cdots \end{pmatrix},
\]
whose first coordinate is non-zero.
\end{proof}
\appendix
\section{Symmetries of~(\ref{LLL})}
For $\alpha \in \C$, define the magnetic translation
\begin{equation*}
R_{\alpha}: u (z) \mto u(z+\alpha) \EE\bigl((\overline z \alpha - z \overline{\alpha})/2\bigr).
\end{equation*}
Let us recall some material from \cite[App.\,A]{SchTho}. We~have
\begin{equation*}
R_{\alpha}u=e^{i(\alpha \cdot \Gamma)}u,
\end{equation*}
with $\alpha=\alpha_1+i\alpha_2$ and $\alpha \cdot \Gamma:= \alpha_1 \Gamma_1 +\alpha_2 \Gamma_2$, where $\Gamma_1$ and $\Gamma_2$ are defined by
\begin{equation}\label{infinitesimal}
\Gamma_1=i(z -\partial_z-\sfrac{\ov z}2), \qquad \Gamma_2=(z +\partial_z+\sfrac{\ov z}2).
\end{equation}
Notice that the operators $R_{\alpha}$ do not commute in general. Indeed, for all $\alpha, \beta \in \C$
\begin{equation}\label{commut}
R_{\alpha} R_\beta = e^{ {\ov \alpha} \beta-\alpha \ov{\beta} } R_{\beta} R_{\alpha}.
\end{equation}
We also record the following formulas
\begin{equation}
\label{propRalpha}
\begin{split}
R_\alpha R_\beta &= R_{\alpha + \beta} \qquad \text{if $\alpha,\beta$ are colinear}, \\
(R_\alpha)^{-1} &= R_{-\alpha} \\
R_\alpha [ u(-z) ] &= [R_{-\alpha} u](-z).
\end{split}
\end{equation}
Furthermore,\vspace*{-3pt}
\begin{align}\label{com0}
R_\alpha \Pi u &= \Pi R_\alpha u \qquad \text{if $|u(x)| \lesssim \langle x \rangle^M$} \\
(R_\alpha)^* &= R_{-\alpha} \qquad \text{on $L^2$}, \nonumber
\end{align}
\vspace*{-1.5\baselineskip}
\par\noindent
and\vspace*{-3pt}
\begin{equation}
\label{r-1}
R_{-\beta} (\alpha \cdot \Gamma) R_\beta =(\alpha \cdot \Gamma) -2 \Im(\alpha \ov{\beta}).
\end{equation}

\section{Proof of Proposition \ref{prop41}}

\label{App1}
Let $u \in \mathcal{E}_{\tau,\gamma}$, and find first a fundamental cell $Q$ with respect to $\mathcal{L}_{\tau, \gamma}$ such that $u$ does not vanish on $\partial Q$. Let $\{ z_k \}_{1\leq k\leq N}$ be the zeroes of $u$ on $Q$, and write\vspace*{-3pt}
\[
\dis u(z) = \EE\bigl(-|z|^2/2 \bigr) \varphi(z) \prod_{k=1}^N \Theta_{\tau}\Bigl(\frac{1}{\gamma}(z-z_{k})\Bigr).
\]
By construction, $\varphi$ does not vanish on $\mathbb{C}$, and is entire; thus, it can be written $\varphi = \EE\bigl(\Psi\bigr)$, with $\Psi$ entire. Furthermore, denoting $A = \mathbb{C} \setminus \bigcup_{a \in \mathcal{L}_{\tau,\gamma}} B(a,\eps)$ for $\eps>0$ small enough, it is easy to see that $|\Theta_\tau(z)| \geq C \EE\bigl({-C|z|^2}\bigr)$ on~$A$. Since $u$ is bounded on~$\mathbb{C}$, this implies that $\Re \Psi (z) \leq C |z|^2$ on~$A$, hence on $\mathbb{C}$ by the maximum principle; but the Borel-Carathéodory theorem implies then that $\Psi$ is actually a polynomial of degree~2. Therefore, we~can write\vspace*{-3pt}
\begin{equation}\label{eqA1}
\dis u(z) = \lambda \EE\bigl({-|z|^2/2 + \alpha z^2 + \beta z}\bigr) \prod_{k=1}^N\Theta_{\tau}\Bigl(\frac{1}{\gamma}(z-z_{k})\Bigr),
\end{equation}
where $\alpha,\beta,\lambda \in \mathbb{C}$. We~take for simplicity $\lambda =1$ in the following. The first periodicity condition of~$\mathcal{E}_{\tau,\gamma}$ requires that $u(z+\gamma) = \EE\bigl({\gamma}(z-\overline{z})/2\bigr) u(z)$. Given the above form of $u$ and~\eqref{star}, this is equivalent to\vspace*{-3pt}
\begin{multline*}
(-1)^N \EE\Bigl(-\frac{1}{2}|z|^2 - \frac{\gamma}{2} \ov{z} + \alpha z^2 + z\bpa{-\frac{\gamma}{2} + 2 \alpha \gamma+ \beta} + \bpa{\alpha \gamma+ \beta - \frac{\gamma}{2}}\gamma \Bigr) \\
= \EE\Bigl(\frac{\gamma}{2}(z-\ov{z})\Bigr) \EE\Bigl({-\frac{1}{2}|z|^2+\alpha z^2 + \beta z}\Bigr).
\end{multline*}
Identifying the coefficients of the polynomials in the exponential, we~see that this equality holds if and only if $\alpha = 1/2$ and $\beta = -{iN\pi}/{\gamma} + {2ik\pi}/{\gamma}$, with $k \in \mathbb{Z}$. The second periodicity condition of $\mathcal{E}_{\tau,\gamma}$ demands that $\dis u(z+\gamma\tau) = \EE\bigl({{\gamma}(\ov{\tau}z-\tau \ov{z})/2}\bigr)u(z)$. For~$u$ as above, and taking~\eqref{star} into account, which yields\vspace*{-3pt}
\[
\Theta_{\tau}\Bigl(\frac{1}{\gamma}(z-z_{k}+\gamma \tau)\Bigr) =\Theta_{\tau}\Bigl(\frac{1}{\gamma}(z-z_{k})+ \tau\Bigr) =- e^{-i \pi \tau}e^{- {2i\pi} (z-z_k)/ {\gamma}}\Theta_{\tau}\Bigl(\frac{1}{\gamma}(z-z_{k})\Bigr),
\]
the periodicity condition is equivalent to\vspace*{-5pt}
\begin{multline*}
\EE\Bigl(-\frac{1}{2}|z|^2 + \frac{1}{2}z^2 -\frac{\gamma}{2} \tau \ov{z} + z(-\frac{\gamma}{2} \ov{\tau} - \frac{2iN\pi}{\gamma} + \gamma\tau + \beta) \Bigr)\\[-3pt]
\times \EE\Bigl((-\frac{\gamma^2}{2} |\tau|^2 + Ni\pi - Ni\pi \tau + \frac{2i\pi}{\gamma} S_N + \frac{\gamma^2}{2} \tau^2 + \beta \gamma \tau)\Bigr)\\[-3pt]
= \EE\Bigl(\frac{\gamma}{2}(\ov{\tau} z - \tau \ov{z})\Bigr) \EE\Bigl(-\frac{1}{2} |z|^2 + \frac{1}{2} z^2 + \beta z\Bigr),
\end{multline*}
where $S_N=\sum_{k=1}^Nz_k$. Identifying the coefficients in the polynomials, this equality holds if and only if $\tau_2 = {N\pi}/{\gamma^2}$ (this comes from the identification of the factor $z$) and
\[
\frac{\gamma^2}{2} \tau^2 - \frac{\gamma^2}{2} |\tau|^2 + Ni\pi - Ni\pi\tau + \frac{2i\pi}{\gamma} S_N + \beta \gamma\tau = 2i \ell \pi, \quad \text{with $\ell \in \mathbb{Z}$},
\]
coming from the constant term. Since $\tau_2 = {N\pi}/{\gamma^2}$, this simplifies to $N + {2}S_N/{\gamma} - N\tau + 2k\tau = 2\ell$, with $\ell \in \mathbb{Z}$. Therefore, $S_N = {\gamma N}(\tau-1)/2+\gamma \ell-\gamma \tau k$. Overall, we~found that~$u$ reads
\[
u(z) = \EE\Bigl({-\frac{1}{2}|z|^2 + \frac{1}{2} z^2 + \frac{i\pi z}{\gamma}(-N + 2k)}\Bigr) \prod_{k=1}^N \Theta_{\tau}\Bigl(\frac{1}{\gamma}(z-z_{k})\Bigr),
\]
with
\[
\sum_{k=1}^N z_k = \frac{\gamma N}{2}(\tau-1)+\gamma \ell-\gamma \tau k. 
\]
There remains to show that $\mathcal{E}_{\tau,\gamma}$ is a vector space of dimension $N$. This follows from two observations: on the one hand, it is a vector space by definition, and on the other hand the number of free parameters in the formula for $u$ is $N$ (the multiplicative constant, and the $N$ zeros which have a prescribed sum).

\section{Useful formulas}

\subsection{Gaussian integral} If $a>0$, $b \in \mathbb{R}$,
\begin{equation}
\label{GaussianIntegral}
\int_{\R} \EE\bigl({-at^2+bt}\bigr) dt = \sqrt{\frac{\pi}{a}} \EE\Bigl(\frac{b^2}{4a}\Bigr).
\end{equation}

\subsection{Poisson summation}
Let us recall the Poisson summation formula (see for instance \cite[Chap.\,VII]{SteinWeiss}): for $g \in \mathscr{S}(\R)$ a 1-periodic function set
\[
\hat{g}(n)=\int_0^1g(x)e^{-2inx}dx,
\]
then for all $ x \in \R$
\begin{equation}\label{PF}
\sum_{n \in \Z}g(x+n)=\sum_{n \in \Z}\hat{g}(n)e^{2i \pi n x}.
\end{equation}
In particular, for $\alpha >0$ and $g: x \mto \EE\bigl({-\alpha x^2}\bigr)$, and by an analytical extension, we~get that for all $z \in \C$
\begin{equation}\label{PF1}
\sum_{n \in \Z} \EE\bigl({-\alpha (z+n)^2}\bigr) =\sqrt{\frac{\pi}{\alpha}}\, \sum_{n \in \Z} \EE\bigl({-{\pi^2 n^2 }/{\alpha}+2 i \pi n z}\bigr).
\end{equation}
\begin{thebibliography}{BGR23}
\bibitem[14]{Cha} K. Chandrasekharan – Elliptic functions, Grundlehren Math. Wissen., vol. 281, Springer-Verlag, Berlin, 1985.
\bibitem[26]{Godement} R. Godement – Analysis. IV, Universitext, Springer, Cham, 2015.
\bibitem[31]{Mumford} D. Mumford – Tata lectures on theta. I, Progress in Math., vol. 28, Birkhäuser Boston, Inc., Boston, MA, 1983.
\bibitem[33]{Nier2} F. Nier – “À propos des fonctions thêta et des réseaux d’Abrikosov”, in Séminaire équations aux dériées partielles 2006-2007, École Polytechnique, Palaiseau, 2007, Exp. no. XII, 27 p.
\bibitem[7]{AftaSerfa} A. Aftalion \& S. Serfaty – “Lowest Landau level approach in superconductivity for the Abrikosov lattice close to Hc2 ”, Selecta Math. (N.S.) 13 (2007), no. 2, p. 183–202.
\bibitem[6]{ABN} A. Aftalion, X. Blanc \& F. Nier – “Lowest Landau level functional and Bargmann spaces for Bose-Einstein condensates”, J. Funct. Anal. 241 (2006), no. 2, p. 661–702.
\bibitem[21]{GGT} P. Gérard, P. Germain \& L. Thomann – “On the cubic lowest Landau level equation”, Arch. Rational Mech. Anal. 231 (2019), no. 2, p. 1073–1128.
\bibitem[42]{SchTho} V. Schwinte \& L. Thomann – “Growth of Sobolev norms for coupled lowest Landau level equations”, Pure Appl. Anal. 3 (2021), no. 1, p. 189–222.
\bibitem[48]{SteinWeiss} E. M. Stein \& G. Weiss – Introduction to Fourier analysis on Euclidean spaces, Princeton Math. Series, vol. 32, Princeton University Press, Princeton, NJ, 1971.
\end{thebibliography}
\end{document}
